% Level 3, Week 5 -- Sorting, Stacks and Queues
\lesson{5}{Sorting, Stacks and Queues}{Sorting puts data in order; a stack remembers what to undo; a queue makes sure everyone gets a turn.}
\lessoninfo{Concepts: selection sort, bubble sort, sorted(key=...), stack, queue, collections.deque}{Tools: Thonny, turtle, git}{CSTA: 3A-AP-15, 3A-AP-14}

\begin{objectives}
  \item The student can trace and code bubble sort and compare it with selection sort by counting swaps. (L3.G2)
  \item The student can sort a list of \texttt{(nickname, points)} tuples with \texttt{sorted(key=..., reverse=True)} to make a leaderboard. (L3.G3)
  \item The student can use a stack (last in, first out) to undo drawing steps and a queue (first in, first out) to take turns. (L3.G3)
  \item The student can suggest ways to make a leaderboard fair and kind. (L3.G9)
\end{objectives}

\begin{materials}
  \item Print: workbook chapter 5 and homework 5 (one per student).
  \item Unplugged: 8 playing cards or paper cards numbered 1--8; 3 plates, books or cups to stack; 4 small toys or pencils for a \enquote{lunch queue}.
  \item On each Mac: Thonny with \texttt{3-advanced/exercises/week05} open; Terminal in \texttt{Documents/git-quest}.
\end{materials}

\section{At a glance}
\begin{timeplan}
  \timeslot{0--8}{Warm-up}{Unplugged human sorting race: selection rules versus bubble rules, counting swaps.}
  \timeslot{8--18}{Concept}{How the two sorts work; \texttt{sorted(key=...)}; plate stack and lunch queue.}
  \timeslot{18--28}{Guided demo}{Bubble sort trace on 4 cards, then the code. Bug hunt: the swap that loses items.}
  \timeslot{28--52}{Hands-on build}{Pairs: Exercise 5.1 Sort Race, 5.3 Undo Paint (stack); fast pairs also 5.2 Leaderboard. Swap driver every 8 minutes. Last 5 minutes: free drawing with undo.}
  \timeslot{52--60}{Wrap-up}{Show-and-tell drawings, fair leaderboard talk, exit ticket, homework.}
\end{timeplan}

\section{Warm-up (0--8 min)}
\begin{unplugged}[Human sorting race]
Shuffle the 8 cards and lay them face up in a row. Copy the row on paper so you can reset it.
\textbf{Selection rules:} find the smallest card, swap it with the first card; then find the smallest of the rest, swap it into second place, and so on.
\textbf{Bubble rules:} start at the left; compare two neighbors; if the left one is bigger, swap them; move one step right; at the end, start again from the left until a whole pass has no swaps.
One student sorts with selection rules while another keeps a tally of \textbf{swaps}. Reset the row; swap jobs and sort with bubble rules. (With 1 student: you keep the tally. With 3--4: two sorters race on two copies of the row.)
Ask: \enquote{Which rules needed more swaps? Which felt easier to follow?} (Bubble usually needs many more swaps; selection needs at most 7.)
\end{unplugged}

\section{Concept (8--18 min)}
\begin{teacherscript}
\enquote{Computers sort all the time: high scores, contacts, songs. Both of today's sorts give the same answer, but they do different amounts of work. \textbf{Selection sort} finds the smallest item and swaps it into place: few swaps, lots of looking. \textbf{Bubble sort} compares neighbors and swaps them: big items \enquote{bubble} to the end, like bubbles rising in a drink.

Python also has a built-in sort: \texttt{sorted(scores, key=points\_of, reverse=True)}. The \textbf{key} is a function that tells Python \emph{what to sort by}, for example the points in each \texttt{("Ada", 120)} pair.}
\end{teacherscript}
Now stack the 3 plates: \enquote{Which plate can you take?} (Only the top one: the last one put on.) Line up the 4 toys: \enquote{Who gets lunch first?} (The one at the front: the first to arrive.)

\begin{analogy}[Plates and lunch lines]
A \textbf{stack} is a pile of plates: you \emph{push} a plate on top and \emph{pop} the top plate off. Last in, first out. That is exactly how \enquote{undo} works: the last thing you did is the first thing undone.
A \textbf{queue} is a lunch line: you join at the back and leave from the front. First in, first out. That is how fair turn-taking works.
\end{analogy}

\section{Guided demo (18--28 min)}
\begin{workedexample}[Bubble sort on 4 cards]
Lay out 5, 2, 4, 1. A student does the bubble rules while you write each row:
\begin{center}
\begin{tabular}{@{}lll@{}}
\toprule
\textbf{Pass} & \textbf{Compare and swap} & \textbf{Row after the pass} \\ \midrule
1 & 5/2 swap, 5/4 swap, 5/1 swap & 2, 4, 1, \textbf{5} \\
2 & 2/4 keep, 4/1 swap & 2, 1, \textbf{4, 5} \\
3 & 2/1 swap & \textbf{1, 2, 4, 5} (5 swaps in total) \\
\bottomrule
\end{tabular}
\end{center}
After every pass, the biggest unsorted card has reached its place (bold), so each pass can stop one card earlier. Now read the code; the inner \texttt{for} loop is one pass:
\pyfile[firstline=22,lastline=32]{../exercises/week05/ex-sort-race/solution.py}
Point at the swap line: \texttt{result[i], result[i + 1] = result[i + 1], result[i]}. Python swaps both at once.
\end{workedexample}

\begin{bughunt}[The swap that loses items]
Show this \enquote{swap} and ask what happens to the cards \texttt{[5, 2]}:
\begin{pycode}
if result[i] > result[i + 1]:
    result[i] = result[i + 1]       # now both are 2!
    result[i + 1] = result[i]       # still 2: the 5 is lost
\end{pycode}
The list becomes \texttt{[2, 2]}: a sort that \emph{loses items}. Act it out with two cards: put the 5 down on top of the 2 \enquote{spot} and the 2 is gone. Fix: the tuple swap from the demo (or a helper variable \texttt{temp}). Then ask: \enquote{What if the comparison were \texttt{<} instead of \texttt{>}?} (It sorts from biggest to smallest: a wrong comparison.)
\end{bughunt}

\section{Hands-on build (28--52 min)}
\begin{pairroles}[Driver and navigator]
\textbf{Driver} types; \textbf{navigator} keeps the cards on the table and acts out what the code does when the pair is unsure. Swap at the 8-minute timer.
With \textbf{1 student}: you are the navigator holding the cards. With \textbf{3}: add a \textbf{tester} who runs \texttt{check\_my\_work.py} and reads the hints; rotate. With \textbf{4}: two pairs.
\end{pairroles}
\begin{enumerate}[itemsep=2pt]
  \item \textbf{Exercise 5.1 Sort Race} (about 8 minutes): finish \texttt{bubble\_sort}. Run it: same sorted cards, but 5 swaps for selection and 16 for bubble.
  \item \textbf{Exercise 5.3 Undo Paint} (about 10 minutes): finish the \texttt{Stack} class. The drawing program is ready: arrow keys draw, \texttt{c} changes color, \texttt{u} undoes. Click the turtle window before pressing keys.
  \item \textbf{Exercise 5.2 Leaderboard} for fast pairs: \texttt{sorted} with a \texttt{key} function, and a fair rule for ties.
\end{enumerate}

\begin{freebuild}[Undo art]
Last 5 minutes: draw something with Undo Paint (a robot, a maze, your initials in blocks) and use \texttt{u} to fix mistakes. Ideas for coders: change \texttt{STEP}, add diagonal keys, or add a key that clears everything by popping until the stack is empty. Commit the code in \texttt{git-quest}.
\end{freebuild}

\section{Wrap-up (52--60 min)}
Show-and-tell: each pair shows a drawing and explains what the stack holds (one stroke per item).

\begin{safety}[Leaderboards: fair and kind]
Show the output of Exercise 5.2 (or read it from the README): Leo and Omar both have 95 and share 3rd place. Ask: \enquote{Is that fair? How would you feel if you were always at the bottom of a board everyone can see?} Ideas from good game design: show nicknames only, never real names; give ties the same place; show personal bests and \enquote{most improved}, not only the winners; let players choose not to appear. A leaderboard should make people want to play again, not feel bad.
\end{safety}

\begin{formativecheck}[Exit ticket]
\textbf{Look for:} the tuple swap; \texttt{range(end)} so \texttt{i + 1} never goes past the end; \texttt{push} using \texttt{append} and \texttt{pop} using \texttt{pop()}.\\
\textbf{Ask:} \enquote{A stack: push A, push B, push C, then pop. What comes out? A queue: A, B, C join the line, then one takes a turn. Who?}\\
\textbf{Answer:} the stack gives C (last in, first out); the queue gives A (first in, first out).
\end{formativecheck}

\begin{reflection}[Talk or write]
\begin{itemize}[itemsep=1pt]
  \item Name one place where you see a stack and one where you see a queue in everyday life or in apps.
  \item Selection sort used fewer swaps. Is fewer swaps always \enquote{faster}? What else does it do a lot of?
\end{itemize}
\end{reflection}

\section{Notes for the teacher}
\begin{commonmistakes}
  \item \textbf{A one-way \enquote{swap}} (\texttt{a = b} then \texttt{b = a}) loses an item. Use \texttt{a, b = b, a}.
  \item \textbf{\texttt{for i in range(len(result))}} with \texttt{result[i + 1]}: \texttt{IndexError} on the last item. Stop one earlier.
  \item \textbf{\texttt{scores = scores.sort()}}: \texttt{.sort()} sorts in place and returns \texttt{None}. Use \texttt{sorted(scores)} for a new list.
  \item \textbf{\texttt{key=points\_of()}} with brackets calls the function too early. Pass the function's name: \texttt{key=points\_of}.
  \item \textbf{\texttt{pop(0)} on a stack} takes the bottom plate. A stack pops from the end: \texttt{self.items.pop()}.
  \item \textbf{Keys do nothing in Undo Paint}: click on the turtle window first so it has the keyboard focus.
\end{commonmistakes}

\begin{differentiation}
\supportpath Do the 4-card bubble sort trace again on paper, then Exercise 5.1 only, then 5.3 with the plate stack beside the keyboard (push = put a plate on, pop = take one off).
\extensionpath Do 5.2 and its \texttt{most\_improved} extra, then add redo to Undo Paint with a second stack.
\end{differentiation}

\begin{stretchchallenge}[Stop early]
If a whole pass of bubble sort makes no swaps, the list is already sorted. Add a variable \texttt{swapped = False} to each pass and stop the outer loop when it stays \texttt{False}. Test it on \texttt{[1, 2, 3, 4, 5]}: how many passes does it need now?
\end{stretchchallenge}

\begin{teachernote}[If you are short on time]
Skip the plates and lunch line (the analogy box is enough). Do 5.1 and 5.3. The Leaderboard and the fair-leaderboard talk can open next week's review session; the queue is practiced in homework 5.
\end{teachernote}
